\section{Foundations of Legal Artificial Intelligence (1991--2016)}
\label{sec:foundations}

\subsection{From legal knowledge to computational reasoning}

The earliest generation of LegalAI was shaped by a fundamental question: can a computer represent
the knowledge and reasoning that legal professionals use to resolve disputes? At the time, the
dominant challenge was not generating natural-sounding legal answers. It was expressing legal
rules, facts, relationships, and arguments in a form that a computational system could process.

One influential approach was case-based reasoning, which draws on the legal practice of examining
previous decisions when addressing a new dispute. Ashley's work on HYPO demonstrated how a system
could compare a fact pattern with earlier cases, identify legally relevant similarities and
differences, and construct arguments using those comparisons \citep{ashley1991}. Rather than
treating a past decision as a simple answer to a new question, this approach made the reasoning
process itself part of the computational problem.

Consider a simplified example. Suppose a court must decide whether a party breached a contractual
obligation. A case-based system might retrieve earlier disputes involving similar obligations,
compare the parties' conduct, and identify facts that strengthen or weaken an argument. If the
current dispute differs from a precedent in a legally significant way, the system must account for
that difference rather than blindly reuse the earlier outcome. This illustrates both the strength
and the difficulty of case-based reasoning: legal relevance depends on more than textual
similarity.

A related approach was the rule-based expert system. Instead of deriving guidance primarily from
comparisons with previous cases, a rule-based system represents knowledge as explicit rules and
applies them to known facts. In simplified form, a rule might state that if specified conditions
are satisfied, a particular legal consequence follows, subject to any applicable exceptions. Such
systems can make their reasoning more explicit because the rules used to reach a conclusion are
represented directly.

However, legal rules rarely operate in isolation. Their application may depend on definitions,
exceptions, procedural requirements, competing authorities, and facts that remain uncertain. A
system that represents a general rule but omits its exceptions can produce an incorrect conclusion
even when its rule-matching procedure works exactly as designed. Similarly, a manually encoded
system requires ongoing maintenance as legislation changes and courts interpret existing provisions
in new ways.

These early approaches therefore established a lasting distinction. Legal reasoning requires both a
representation of legal knowledge and a method for applying that knowledge to a particular dispute.
A system may possess explicit rules but fail to identify the relevant facts; it may retrieve a
similar case but fail to explain why the differences matter. This distinction remains important in
modern legal AI, where retrieval and generation are often treated as separate components of a
larger reasoning process.

\subsection{Rule-based systems and legal knowledge representation}

Rule-based expert systems sought to make legal knowledge operational. Legal experts or researchers
encoded domain knowledge into rules, facts, and decision procedures, allowing the system to apply
that knowledge to particular inputs. This approach was attractive where a legal task could be
described through a relatively stable set of conditions and consequences.

The main advantage was control. Unlike a statistical model whose internal representations may be
difficult to interpret, an explicitly encoded rule can often be inspected and traced. A developer
can identify which condition was satisfied and which rule contributed to the outcome. This makes
rule-based approaches conceptually useful for tasks requiring transparent, predefined logic.

Their limitations become clearer when legal knowledge grows in size and complexity. Rules may
interact, exceptions may override general conditions, and different authorities may apply in
different circumstances. Manually encoding these relationships can become expensive, while
incomplete or outdated rules can lead to systematic errors. The challenge is not merely writing
more rules; it is maintaining a consistent representation of the legal system and the relationships
among its components.

This led to broader work on legal knowledge representation, including ontologies and formal models
of legal concepts. A legal ontology defines concepts within a domain and specifies relationships
among them. For example, it may distinguish a statute from a judicial decision, represent a legal
provision as belonging to a particular enactment, or connect a legal obligation with the
circumstances under which it applies. Such representations can support structured querying,
information integration, and reasoning over explicit relationships.

Ontologies and knowledge graphs are related but not identical. An ontology defines a conceptual
vocabulary and the types of relationships that may exist; a knowledge graph uses entities and
relationships to represent particular information. A graph might connect a statutory provision to an
amendment, a judgment that interprets it, and a legal issue addressed by that judgment. The value of
these connections depends on their accuracy, completeness, and provenance.

Nevertheless, structured knowledge representation also has costs. Legal concepts can be contested,
the meaning of a provision may change through judicial interpretation, and relationships may require
human expertise to establish reliably. Formal representations are therefore useful but do not remove
the need to interpret legal texts and verify the authority of the underlying information.

\subsection{Early legal information retrieval}

A different but complementary line of research focused on finding legal documents relevant to a
user's information need. Legal information retrieval is not identical to legal reasoning. Retrieval
aims to identify and rank potentially relevant documents or passages, while reasoning determines how
the retrieved material should be interpreted and applied.

Traditional information retrieval methods relied heavily on lexical evidence: the words appearing in
a query and the documents being searched. This is valuable in law because exact wording, statutory
section numbers, party names, and established legal phrases can be decisive. A query that mentions a
particular provision may require that exact authority rather than a passage that merely discusses a
similar concept.

However, lexical matching can miss relevant material when the query and the authoritative document
use different expressions. A user may describe a legal problem in everyday language, while the
relevant judgment expresses the same issue using specialized terminology. Conversely, two documents
may share many words while addressing different legal issues. This creates a recurring challenge: the
words used to describe a legal problem are not always the words used in the authority that resolves
it.

Case retrieval adds another complication. A relevant precedent is not necessarily one that discusses
the same topic. Its value may depend on the similarity of the legally material facts, the legal
issue, the reasoning, and the outcome. Research on legal case retrieval has consequently developed
specialized datasets and models for identifying precedents relevant to a given case, rather than
relying only on general-purpose document matching \citep{feng2024}.

By the end of this period, the main approaches had established complementary foundations. Rule-based
systems offered explicit knowledge and controlled inference; case-based reasoning emphasized
comparison with prior disputes; ontologies represented legal concepts and their relationships; and
information retrieval provided mechanisms for finding relevant legal material. Each addressed part of
the problem, but none alone solved the full challenge of locating authoritative evidence, preserving
its context, and producing a reliable answer.

\subsection{Lessons from the foundational period}

The importance of this period lies not in a single winning architecture, but in the problems it
exposed. Legal knowledge must be represented carefully, relevant authorities must be found, and
reasoning must account for the relationship between a general rule and the facts of a specific case.
These requirements remain relevant even as the underlying computational methods change.
Table~\ref{tab:foundational} summarizes the main foundational approaches.

\begin{table*}[tp]
\centering
\caption{Foundational approaches to LegalAI (1991--2016).}
\label{tab:foundational}
\footnotesize
\setlength{\tabcolsep}{5pt}
\renewcommand{\arraystretch}{1.3}
\begin{tabularx}{\textwidth}{>{\raggedright\arraybackslash}p{2.6cm} >{\raggedright\arraybackslash}X >{\raggedright\arraybackslash}X >{\raggedright\arraybackslash}X}
\toprule
\textbf{Approach} & \textbf{Core idea} & \textbf{Main strength} & \textbf{Principal limitation} \\
\midrule
Case-based reasoning & Compare a current dispute with previous cases and construct arguments from relevant similarities and differences. & Supports precedent-oriented reasoning and case comparison. & Requires identifying legally meaningful similarities and differences; textual similarity alone is insufficient. \\
Rule-based expert systems & Encode legal knowledge as explicit rules and apply them to known facts. & Transparent, inspectable rule application within a defined domain. & Difficult to maintain at scale when rules interact, exceptions arise, and legal authorities change. \\
Legal ontologies & Define legal concepts and their relationships in a structured vocabulary. & Supports consistent representation and structured querying of legal knowledge. & Requires careful modelling and maintenance; formal relationships do not automatically resolve legal ambiguity. \\
Early legal information retrieval & Search and rank legal documents using query terms and document evidence. & Helps users navigate large collections of statutes and judicial decisions. & Lexical matching can miss conceptually relevant documents or retrieve passages without sufficient legal context. \\
Integrated knowledge-based systems & Combine structured knowledge, retrieval, and reasoning components. & Brings different capabilities together within a single application. & System quality depends on the completeness, consistency, and coordination of its components. \\
\bottomrule
\end{tabularx}
\end{table*}

The approaches in Table~\ref{tab:foundational} should not be interpreted as mutually exclusive or as
a strict sequence in which one disappeared when another emerged. They represent different ways of
addressing different parts of the legal information problem. Modern systems continue to use lexical
retrieval, explicit legal rules, structured metadata, and precedent comparison, even when their
generators rely on neural language models.

The next phase changed the way these systems acquired and processed linguistic knowledge. Instead of
depending primarily on manually encoded rules or separately engineered features, researchers
increasingly trained models to learn representations from large text collections. This shift made it
possible to build more flexible legal text-processing systems, but it also introduced new questions
about domain adaptation, evaluation, and the reliability of model predictions.
